\subsection{Quasi-invariant measures}

PROPOSITION 11. --- Let G be a locally compact group, \(\mu\) a left Haar measure on G. For a measure \(\nu \neq 0\) on G to be left quasi-invariant, it is necessary and sufficient that \(\nu\) be equivalent to \(\mu\) .

Sufficiency is obvious. Let \(\nu \neq 0\) be a left quasi-invariant measure, and let us show that \(\nu\) is equivalent to \(\mu\) . We can restrict ourselves to the case that \(\nu > 0\) . Let A be a compact subset of G. We will show, as will establish the proposition, that the conditions \(\mu(A) = 0\) , \(\nu(A) = 0\) are equivalent (Ch. V, §5, No.~5, Th. 2).

\begin{enumerate}
\def\labelenumi{\alph{enumi})}
\tightlist
\item
  For every \(f \in \mathcal{K}_{+}(\mathrm{G})\) , the function \((x, y) \mapsto f(x)\varphi_{\mathrm{A}}(xy)\) on \(G \times G\) is \((\nu \otimes \mu)\) -integrable, because it is upper semi-continuous, bounded, and its support is contained in the compact set \(K \times K^{-1}A\) if one sets \(K = Supp f\) . Therefore, by the Lebesgue--Fubini theorem,
\end{enumerate}

\[
\int d \nu (y) \int \varphi_ {\mathrm{A}} (x y) f (x) d \mu (x) = \int f (x) d \mu (x) \int \varphi_ {\mathrm{A}} (x y) d \nu (y).\tag{36}
\]

\begin{enumerate}
\def\labelenumi{\alph{enumi})}
\setcounter{enumi}{1}
\tightlist
\item
  Suppose \(\nu(\mathbf{A}) = 0\) . By hypothesis, \(\nu(x\mathbf{A}) = 0\) for all \(x \in G\) , therefore the right side of (36) is zero. Therefore there exists a \(\nu\) -negligible set \(N_{f}\) such that, for \(y \notin N_{f}\) ,
\end{enumerate}

\[
0 = \int \varphi_ {\mathrm{A}} (x y) f (x) d \mu (x) = \Delta_ {\mathrm{G}} (y) ^ {- 1} \int \varphi_ {\mathrm{A}} (x) f (x y ^ {- 1}) d \mu (x).\tag{37}
\]

Let B be a compact subset of G such that \(\nu(\mathrm{B})\neq0\) , and take for f a function in \(\mathcal{K}_{+}(\mathrm{G})\) equal to 1 on \(AB^{-1}\) . There then exists a \(y\in B\) such that (37) is verified. But since \(\varphi_{\mathrm{A}}(x)f(xy^{-1})=\varphi_{\mathrm{A}}(x)\) for \(y\in B\) , this proves that \(\mu(\mathrm{A})=0\) .

\begin{enumerate}
\def\labelenumi{\alph{enumi})}
\setcounter{enumi}{2}
\tightlist
\item
  Suppose \(\mu(A) = 0\) . Then, for every \(f \in \mathcal{K}_{+}(G)\) , the left side of (36) is zero, hence also the right side. Consequently, there exists a locally \(\mu\) -negligible set \(M\) such that \(\int \varphi_{A}(xy) d\nu(y) = 0\) for \(x \notin M\) . Since \(\mu \neq 0\) , it follows that \(\nu(xA) = 0\) for some \(x \in G\) , whence \(\nu(A) = 0\) .
\end{enumerate}

Applying Prop. 11 to \(G^{0}\) , one sees that the right quasi-invariant measures are identical with the left quasi-invariant measures. They are called simply quasi-invariant measures on G.

